% AUTORIDAD DE FUENTE: Exposé IX, título del §9 y §9.1,
% líneas 2194--2214; índices de escaneo 429--430, folios 418--419.
% Las capas francesas D/C no presentan diferencias materiales.

\medskip
\noindent
9. \underline{Definición del emparejamiento de monodromía}
$\underline{M}_{\ell}\otimes\underline{M}'_{\ell}
 \longrightarrow \mathbb{Z}_{\ell S}$.

\medskip
\noindent
9.1. En esta sección, supongamos que $S$ es henseliano y que $A_K$
tiene reducción semiestable sobre $S$, o, equivalentemente, que $A'_K$
tiene reducción semiestable sobre $S$ (3.5.1). Tenemos entonces sobre
$S$ dos $\mathbb{Z}$-módulos constantes torcidos $\underline{M}$ y
$\underline{M}'$ (5.5.2). Sobre un recubrimiento finito étale
conveniente de $S$ (y por tanto sobre el propio $S$ si $S$ es
estrictamente local, es decir, si $k$ es separablemente cerrado), se
descomponen en grupos constantes de valores $M$ y $M'$,
respectivamente, isomorfos no canónicamente a $\mathbb{Z}^{r}$, donde
$r$ es la dimensión común de los toros maximales de $A_o$ y $A'_o$.
Para todo número primo $\ell$, consideraremos los haces $\ell$-ádicos
correspondientes

\begin{equation}\tag{9.1.1}
\underline{M}_{\ell}=\underline{M}\otimes_{\mathbb{Z}}\mathbb{Z}_{\ell}
 \quad , \qquad
\underline{M}'_{\ell}=\underline{M}'\otimes_{\mathbb{Z}}\mathbb{Z}_{\ell}
 \qquad ,
\end{equation}
cuyo significado geométrico queda explicitado por las fórmulas (5.5.7)
y (5.5.8). En esta sección definiremos un emparejamiento canónico,
llamado el \underline{emparejamiento de monodromía},

\begin{equation}\tag{9.1.2}
u_{\ell}:\underline{M}_{\ell}\otimes\underline{M}'_{\ell}
 \longrightarrow\mathbb{Z}_{\ell} \qquad ,
\end{equation}
que se definirá en términos de los grupos de Barsotti--Tate
pro-$\ell$ $T_{\ell}(A^{o})$ y $T_{\ell}({A'}^{o})$, e incluso en
términos de la fibra genérica $T_{\ell}(A_K)$ del primero de estos
grupos (a reserva, cuando $\ell=p$, de disponer del teorema de Tate
[33, Th.~4], es decir, por ahora, siempre que estemos entonces en
características desiguales). Cuando $\ell\neq p$, esto equivale a decir
que el emparejamiento (9.1.2) se describe en términos del
$\pi$-módulo de Galois

\begin{equation}\tag{9.1.3}
U=T_{\ell}(A_K)(\overline K)=T_{\ell}(A_{\overline K})
\end{equation}
considerado en (2.5.3).

%%%%%%%%%%  scan idx 430  |  folio 419  %%%%%%%%%%

Más precisamente, veremos que conocer el emparejamiento equivale
esencialmente a conocer la acción del grupo de monodromía $I$ sobre
$U$.
